\documentclass[11pt]{article}

\usepackage[margin=1in]{geometry}
\usepackage[T1]{fontenc}
\usepackage[utf8]{inputenc}
\usepackage{lmodern}
\usepackage{booktabs}
\usepackage{longtable}
\usepackage{array}
\usepackage{enumitem}
\usepackage{xcolor}
\usepackage{hyperref}
\usepackage{titlesec}
\usepackage{fancyhdr}
\usepackage{xurl}
\usepackage{array}
\usepackage{longtable}
\usepackage{booktabs}

\hypersetup{
  colorlinks=true,
  linkcolor=blue,
  urlcolor=blue,
  citecolor=blue
}

\setlist[itemize]{leftmargin=1.4em}
\setlist[enumerate]{leftmargin=1.6em}
\renewcommand{\arraystretch}{1.16}
\setlength{\tabcolsep}{4pt}
\titleformat{\section}{\Large\bfseries}{\thesection}{0.75em}{}
\titleformat{\subsection}{\large\bfseries}{\thesubsection}{0.75em}{}

\pagestyle{fancy}
\fancyhf{}
\lhead{E50 Voice Command Machine}
\rhead{Final Project Report}
\cfoot{\thepage}

\newcommand{\code}[1]{\texttt{\detokenize{#1}}}
\newcommand{\statuspass}{\textbf{PASS}}
\newcommand{\statusqualified}{\textbf{QUALIFIED}}
\newcommand{\statuspartial}{\textbf{PARTIAL}}
\newcommand{\statusna}{\textbf{NOT ESTABLISHED}}

\title{\textbf{Final Report E50 Voice Command Machine}\\Loreen Anne P. De Guzman}
\author{MEng Artifical Intelligence, University of the Philippines Diliman}
\date{October 3, 2026}

\begin{document}
\maketitle

\begin{center}
\begin{tabular}{ll}
\toprule
Program name & E50 Voice Command Machine, final frozen Raspberry Pi deployment \\
Final experiment & \code{E50_BG_REVISED_VOCAB_ORIGINAL_PRESERVE_E41_INIT} \\
Target hardware & Raspberry Pi 5 \\
Repository package & \code{GITHUB_PACKAGE_FINAL_20261002} \\
Command model SHA256 & \code{BC8AC64CED4305DDF43BF4DE377F3CC636A764EFF339CD9F92AF06340C50B8FF} \\
Wake model & \code{E37_TARGETED_COLOR_VOLUME_FIX} \\
Threshold policy & \code{E40_NON_LED_COMMAND_GUARDRAIL_THRESHOLDS} \\
\bottomrule
\end{tabular}
\end{center}

\vspace{0.5em}
\noindent
This report is the technical summary of the final E50 evidence. It documents the E50 implementation, model, dataset, thresholds, router, response assets, runtime, GUI boundary, benchmark evidence, repository package, verified Demo Day metrics, benchmark audit, post-freeze GUI clarification, and independent E53 documentation.

\tableofcontents
\newpage

\section{Executive Summary}

E50 is the final Machine Exercise 2 Voice Command Machine delivery: an offline Raspberry Pi 5 voice-command system with a frozen 19-label command vocabulary, a separate wake gate, a command CNN, a confidence/rejection policy, deterministic local routing, local response audio, touchscreen GUI operation, and persistent return-to-listening behavior.

The final frozen recognition/action stack is:

\begin{center}
\code{Pi microphone -> E37 wake -> command capture -> log-Mel preprocessing -> E50 command CNN -> E40 guardrail -> deterministic router -> local action -> local WAV response -> return to listening}
\end{center}

The final physical Raspberry Pi benchmark used 20 trials: 19 valid command-label trials plus one UNKNOWN/no-action trial. Wake success was 19/20 = 95.0\% overall and 18/19 = 94.74\% for valid command trials. Raw command classification was 13/19 = 68.42\%. End-to-end action success was 10/19 = 52.63\%. All accepted/executed valid commands produced the correct action, giving accepted-action precision of 10/10 = 100\%. All non-executed trials were safely rejected without a routed action, giving safe rejection of 10/10 = 100\%. The UNKNOWN/no-action trial was safe at 1/1 = 100\%. The runtime returned to listening after all 20 trials.

These metrics deliberately remain separate. The report does not collapse them into a single ``overall accuracy'' number. Raw classification, confidence acceptance, action execution, safe rejection, and end-to-end success measure different layers of the system.

The final system demonstrates strong action safety, conservative rejection behavior, fast local inference, offline operation after setup, deterministic routing, and reproducible frozen identity. Its main limitation is command-level recognition coverage across the full vocabulary: raw classification was moderate at 13/19, and end-to-end action success was 10/19 in the final physical benchmark.

This report also preserves project boundaries. E53/VCM2 is an independent experiment. It is not an E50 model, training stage, benchmark, deployment component, improvement, or replacement. No E53 implementation file, model, dataset, configuration, or evidence artifact is part of the frozen E50 system.

\section{Project Identity and Final Frozen State}

\begin{longtable}{p{0.31\linewidth}p{0.61\linewidth}}
\toprule
\textbf{Field} & \textbf{Final E50 value} \\
\midrule
Project & ME2 Voice Command Machine, final E50 Raspberry Pi delivery \\
Final experiment ID & \code{E50_BG_REVISED_VOCAB_ORIGINAL_PRESERVE_E41_INIT} \\
Target hardware & Raspberry Pi 5 \\
System type & Offline Raspberry Pi 5 voice-command system for fixed command labels and local actions \\
Wake model & \code{E37_TARGETED_COLOR_VOLUME_FIX} \\
Wake threshold & 0.90 \\
Command model & \code{E50_BG_REVISED_VOCAB_ORIGINAL_PRESERVE_E41_INIT} \\
Command model SHA256 & \code{BC8AC64CED4305DDF43BF4DE377F3CC636A764EFF339CD9F92AF06340C50B8FF} \\
Threshold policy & \code{E40_NON_LED_COMMAND_GUARDRAIL_THRESHOLDS} \\
Command default threshold & 0.90, with documented label-specific threshold behavior \\
Runner & \code{scripts/run_vcm_touchscreen_gui.sh} \\
GPIO & Disabled in final delivery \\
Response & Local WAV/action response enabled \\
Final vocabulary size & 19 labels \\
\bottomrule
\end{longtable}

The final E50 vocabulary is:

\begin{enumerate}
\item \code{PLAY_MUSIC}
\item \code{WEATHER}
\item \code{TIME}
\item \code{LIGHT_ON}
\item \code{LIGHT_OFF}
\item \code{LIGHT_DIM}
\item \code{BRIGHTNESS}
\item \code{TIMER}
\item \code{ALARM}
\item \code{TEMPERATURE}
\item \code{NEXT}
\item \code{PAUSE}
\item \code{STOP}
\item \code{VOLUME_UP}
\item \code{VOLUME_DOWN}
\item \code{CREATE_REMINDER}
\item \code{LIST_REMINDERS}
\item \code{CALL}
\item \code{MESSAGE}
\end{enumerate}

\noindent
\textbf{Important vocabulary boundary:} \code{COLOR} is not a final E50 class. The final E50 system uses \code{LIGHT_DIM}. Any evidence involving \code{COLOR} belongs to an earlier experiment, collective benchmark schema, or historical command-design context unless explicitly identified otherwise.

\section{Deployment Scope and Objectives}

\subsection{Scope}

The deployed E50 system covers:

\begin{itemize}
\item Raspberry Pi 5 runtime and GUI launcher.
\item E37 wake detection.
\item E50 command-recognition CNN.
\item E40 confidence/rejection policy.
\item 19-label command vocabulary.
\item Deterministic router/action layer.
\item Local response WAV assets and audio playback.
\item Offline operation after setup.
\item Validation, benchmark, integrity, reproducibility, and technical documentation.
\end{itemize}

The system is not a cloud assistant, not an LLM, not an online ASR service, and not an arbitrary transcription system. It is a wake-gated fixed-vocabulary command classifier with deterministic local action routing.

\subsection{Objectives}

\begin{enumerate}
\item Build a local Raspberry Pi voice-command machine that can wake, classify commands, route actions, play responses, and return to listening.
\item Preserve the complete final 19-label vocabulary without a runtime whitelist.
\item Demonstrate offline operation without cloud or network inference during verified operation.
\item Validate safety behavior for unsupported, silent, wrong, uncertain, or below-threshold input.
\item Report benchmark metrics by layer: wake, raw classification, acceptance, action correctness, safe rejection, end-to-end success, and runtime recovery.
\item Package the project so a technical reader can inspect requirements, architecture, implementation, engineering history, validation, limitations, reproduction steps, and integrity records.
\end{enumerate}
\newpage
\section{Architecture}

\begin{longtable}{p{0.29\linewidth}p{0.63\linewidth}}
\toprule
\textbf{Stage} & \textbf{Role in E50} \\
\midrule
Pi microphone & Captures physical audio input. \\
E37 wake model & Detects the wake phrase and gates command capture. \\
Command capture & Records the command window after wake acceptance. \\
Log-Mel preprocessing & Converts captured command audio into fixed features. \\
E50 command CNN & Predicts one of the 19 final command labels. \\
E40 policy & Accepts or rejects predictions using frozen confidence thresholds. \\
Router/action layer & Maps accepted labels to deterministic local actions. \\
Local response WAV & Plays the corresponding local response audio. \\
Return to listening & Restores persistent runtime state after each trial. \\
\bottomrule
\end{longtable}

The GUI is an interface and launcher. Recognition and routing are runtime responsibilities downstream of the frozen E37/E50/E40 stack.

\subsection{Recovered E37 Wake-Recording Provenance}

Recovered E37 provenance evidence documents the original project-specific Raspberry Pi \code{Hey Pi} wake-recording resource. The recovered manifest is
\path{/home/loreenanne/vcm_pi_package/pi_validation/wake_validation_hey_pi/manifest.csv}. It contains 99 wake-stage recordings: 50 \code{WAKE} / \code{hey pi} rows, 15 \code{UNKNOWN} rows, 14 \code{COLOR} targeted rows, 14 \code{VOLUME_UP} targeted rows, 3 \code{LIGHT_ON} command-as-wake-probe rows, and 3 \code{PLAY_MUSIC} command-as-wake-probe rows. The manifest records 64 adaptation-designated rows and 35 holdout-designated rows. All rows are documented as 4-second, 16 kHz recordings using input device \code{plughw:2,0}.

This evidence establishes the existence and composition of the E37-related \code{Hey Pi} recording set. It does not independently establish the exact final E37 training-row membership, speaker/recordist count, training epochs, optimizer, validation metrics, augmentation count, or row-level mapping from the 99 recordings to the final E37 checkpoint. The E37 wake-recording lineage is also separate from the collective Gold Dataset lineage used for the E50 command model.

\section{Dataset and Training Provenance}

\subsection{Collective Gold Dataset Lineage}

E50 was developed using selected material from the class collective Gold Dataset:\\
GoogleDrive https://drive.google.com/drive/folders/1k2RklcsU8ull9-4fagl-8zTfecPaQWBZ\\
HuggingFace https://huggingface.co/datasets/airimonda/ai231-me2-voice-commands\\
Documentation https://claude.ai/artifact/PPNHMWd5rx9qcXTdcukV7s
rather than using the entire collective dataset unchanged. The project-specific dataset was further prepared and balanced for the 19-command E50 vocabulary, with training-only augmentation and a separate Pi deployment-adaptation component. The resulting dataset therefore retains lineage to the collective dataset while incorporating project-specific engineering data.

The final E50 training manifest contains 16,100 rows: 13,070 active-project rows, 2,800 selected \code{VCM\_BALANCED} rows, and 230 separately identified E41 Pi adaptation rows. The first two components are collective-derived/project-local command-data branches; the E41 rows are project-specific deployment adaptation data and should not be described as Gold Dataset rows unless additional evidence establishes that link.

\code{VCM\_MASTER} is a project-local Dataset2 / collective-family representation derived from selected collective dataset material; it is not claimed to be the entire Gold Dataset. \code{VCM\_BALANCED} is a training-derived subset of \code{VCM\_MASTER} and therefore inherits its collective-family provenance. E53/VCM2 remains an independent experiment and is not merged into E50.


\subsection{Dataset A: VCM\_MASTER}

\begin{longtable}{p{0.42\linewidth}p{0.46\linewidth}}
\toprule
\textbf{Item} & \textbf{Value} \\
\midrule
Dataset name & \code{VCM_MASTER} \\
Total samples & 36,622 \\
Classes & 16 \\
Training samples & 27,130 \\
Validation samples & 4,734 \\
Test samples & 4,758 \\
Speaker/group identities & 395 \\
Speaker leakage & Zero train/validation/test overlap documented \\
\bottomrule
\end{longtable}

\subsection{Dataset B: VCM\_BALANCED}

\begin{longtable}{p{0.42\linewidth}p{0.46\linewidth}}
\toprule
\textbf{Item} & \textbf{Value} \\
\midrule
Dataset name & \code{VCM_BALANCED} \\
Training samples & 15,268 \\
Original samples & 13,801 \\
Augmented samples & 1,467 \\
Relationship to Dataset A & Derived from Dataset A training material; did not introduce validation/test leakage \\
\bottomrule
\end{longtable}

\subsection{Final E50/BG Training Manifest: 16,100 rows}

\begin{longtable}{p{0.52\linewidth}r}
\toprule
\textbf{Source} & \textbf{Rows} \\
\midrule
\code{active_project_dataset} & 13,070 \\
\code{dataset2_vcm_balanced} & 2,800 \\
\code{e41_reconstructed_adaptation} & 230 \\
\midrule
\end{longtable}
\newpage
The final data-kind composition was:

\begin{longtable}{p{0.52\linewidth}r}
\toprule
\textbf{Kind} & \textbf{Rows} \\
\midrule
Real & 15,260 \\
Augmented & 480 \\
Real multi-sensor & 200 \\
Synthetic & 160 \\
\midrule
\textbf{Total} & \textbf{16,100} \\
\bottomrule
\end{longtable}

The final distribution was controlled rather than perfectly balanced. The engineering rule retained E41 lineage where applicable, capped active-project original data at 700 rows per label, and added up to 200 Dataset2 rows where usable support existed. This yielded 14 labels with 900 rows and five labels with 700 rows. The five 700-row labels were \code{BRIGHTNESS}, \code{CREATE_REMINDER}, \code{LIST_REMINDERS}, \code{CALL}, and \code{MESSAGE}.

This is best described as controlled class balancing with bounded auxiliary-data addition. Synthetic and augmented data are training-domain support, not independent real-speaker deployment evidence.

\subsection{Training History}

The final selected checkpoint was epoch 7, selected by best internal validation loss.

\begin{longtable}{rrrrrl}
\toprule
\textbf{Epoch} & \textbf{Train loss} & \textbf{Train acc.} & \textbf{Val. loss} & \textbf{Val. acc.} & \textbf{Status} \\
\midrule
7 & 0.3458 & 0.8801 & 1.3988 & 0.6197 & Best / selected \\
11 & 0.1946 & 0.9322 & 2.1883 & 0.6018 & Later / not selected \\
\bottomrule
\end{longtable}

The later training trajectory shows a classical overfitting pattern: training loss decreased and training accuracy increased, while validation loss worsened and validation accuracy did not improve meaningfully. However, this does not mean the deployed E50 checkpoint was a late overfit model. Epoch 7 was the selected checkpoint; later overfit epochs were not deployed. Deployment failures should therefore be attributed by evidence-supported layer rather than blamed entirely on overfitting.

\section{Execution Timeline and Engineering Process}

The project followed an evidence-gated engineering process rather than a single uncontrolled training run. Earlier experiments diagnosed model, threshold, router, action, audio-output, and deployment issues separately. Candidate models and recovery experiments were rejected when they introduced regressions or lacked adequate independent validation.

The final engineering process preserved the following principles:

\begin{itemize}
\item keep model development separate from final validation;
\item keep wake detection, command classification, confidence policy, routing, action execution, response playback, and GUI behavior separable in evidence;
\item avoid post-hoc threshold tuning against final benchmark evidence;
\item preserve all 19 final labels as callable/testable;
\item report demo readiness separately from runtime callability;
\item document limitations instead of patching the frozen system after benchmark closure.
\end{itemize}

The project record documents major September 30 closure activities: offline operation evidence, safety sweep evidence, Pi performance measurement, all-19 final benchmark evidence, and final artifact integrity work. This report presents that project evidence together with the clarified timing, benchmark-audit, GUI-boundary, and E53-separation language used in the final submission package.

\section{Final Benchmark Study}

\subsection{Benchmark Population}

The final physical Raspberry Pi benchmark population was:

\begin{itemize}
\item 19 valid command-label trials, one for each final E50 label;
\item 1 UNKNOWN/no-action trial;
\item 20 total physical Pi 5 trials.
\end{itemize}

The benchmark recorded wake detection, raw predicted label, command confidence, threshold, acceptance, expected action, actual action, action correctness, response playback, return-to-listening, end-to-end success, and failure layer.

\subsection{Final Pi Results}

\begin{longtable}{p{0.35\linewidth}p{0.22\linewidth}p{0.35\linewidth}}
\toprule
\textbf{Metric} & \textbf{Result} & \textbf{Meaning} \\
\midrule
Wake success & 19/20 = 95.0\% & Wake accepted across all benchmark trials. \\
Wake success on valid commands & 18/19 = 94.74\% & Wake accepted for valid command-label trials. \\
Raw command classification & 13/19 = 68.42\% & Correct raw predicted command label among valid commands. \\
Accepted valid commands & 10/19 = 52.63\% & Valid commands accepted by E40 policy. \\
Accepted-correct & 10 & Accepted command produced correct corresponding action. \\
Accepted-wrong & 0 & No accepted command produced a wrong action. \\
Accepted-action precision & 10/10 = 100\% & Correct action among accepted/executed valid commands. \\
Safe rejection among non-executed & 10/10 = 100\% & Non-executed trials produced no routed action. \\
End-to-end action success & 10/19 = 52.63\% & Valid command trial completing the defined successful E50 action path, including correct command outcome, accepted execution, response handling, and return-to-listening. \\
UNKNOWN/no-action safety & 1/1 = 100\% & Unsupported phrase produced no action. \\
Return-to-listening & 20/20 = 100\% & Runtime returned to listening after every trial. \\
\bottomrule
\end{longtable}

The distinction matters. Raw command classification measures the CNN label. Acceptance measures the E40 policy. Accepted-action precision measures whether accepted commands routed to the correct action. Safe rejection measures conservative no-action behavior. End-to-end action success requires the complete successful path.

\subsection{Observed Failure Layers}

The final Item 16 per-trial data show that unsuccessful valid-command trials were not downstream action-routing bugs. They included wrong raw predictions safely rejected, correct raw predictions rejected below threshold, and one wake rejection. Accepted wrong actions were zero. This supports the conclusion that the final benchmark failures were recognition, wake, or confidence/rejection failures rather than erroneous executed actions.

\section{Performance and Efficiency}

\subsection{Model Size and Computational Cost}

\begin{longtable}{p{0.42\linewidth}p{0.23\linewidth}p{0.25\linewidth}}
\toprule
\textbf{Metric} & \textbf{Value} & \textbf{Definition} \\
\midrule
Parameters & 66,483 & Final E50 command CNN parameters. \\
Weights & 251,734 bytes & Final E50 weights artifact size. \\
Weights + normalization & 252,464 bytes & Weights plus normalization artifact. \\
MACs & 85,774,144 & Conv2D/Dense MACs per 4-second input. \\
Pi model load & 1.8566 s & Model load time in Item 15 benchmark. \\
\bottomrule
\end{longtable}
\newpage
\subsection{CNN Inference Latency and CNN-Only RTF}

\begin{longtable}{p{0.42\linewidth}p{0.23\linewidth}p{0.25\linewidth}}
\toprule
\textbf{Measurement} & \textbf{Value} & \textbf{Boundary} \\
\midrule
Mean CNN inference & 31.82 ms & CNN inference only. \\
P50 CNN inference & 31.39 ms & CNN inference only. \\
P95 CNN inference & 34.40 ms & CNN inference only. \\
Maximum CNN inference & 38.31 ms & CNN inference only. \\
CNN-only RTF mean & 0.0080 & CNN inference time / 4-second command input. \\
CNN-only RTF P95 & 0.0086 & CNN inference time / 4-second command input. \\
\bottomrule
\end{longtable}

CNN-only RTF is not full end-to-end RTF, wake-to-action RTF, wake-to-response RTF, or acoustic response latency.

\subsection{Post-Closure Response Playback-Start Timing}

A post-closure passive timing audit measured local response playback-start latency without modifying the frozen runtime. The measurement population was five accepted/executed trials. The boundary was:

\begin{center}
\code{command WAV capture completion -> first observed local aplay response-WAV process}
\end{center}

\begin{longtable}{p{0.42\linewidth}p{0.23\linewidth}p{0.25\linewidth}}
\toprule
\textbf{Measurement} & \textbf{Value} & \textbf{Qualification} \\
\midrule
N & 5 accepted/executed trials & Qualified post-closure passive audit. \\
Mean playback-start latency & 66.172 ms & Not acoustic speaker-onset. \\
P50 playback-start latency & 64.837 ms & Not full microphone-to-speaker latency. \\
P95 playback-start latency & 73.999 ms & Not full wake-to-response latency. \\
P99 playback-start latency & 73.999 ms & Same measurement boundary. \\
Maximum playback-start latency & 73.999 ms & Same measurement boundary. \\
Command-pipeline RTF mean & 0.016543 & Capture completion to response playback-start / 4 s. \\
Command-pipeline RTF P95 & 0.018500 & Capture completion to response playback-start / 4 s. \\
\bottomrule
\end{longtable}

This measurement strengthens the timing evidence, but it does not establish acoustic speaker-onset latency, full microphone-to-speaker latency, full wake-to-action latency, or full wake-to-response latency.

\section{Safety, FAR, and Failure Analysis}

\subsection{Bounded FAR}

The established false-accept result is:

\begin{center}
\textbf{Bounded command-window FAR: 0/4 = 0.0\%}
\end{center}

The definition is: false accepts among valid out-of-scope command-window safety prompts, where a false accept means the frozen E50 system accepted and executed a command.

This is a bounded command-window FAR. It is not a broad environmental false-accept rate, universal false-wake rate, or proof of robustness to arbitrary background audio. Broad environmental FAR was not established by the available evidence.

\subsection{Safety Findings}

The final benchmark supports these safety conclusions within the tested population:

\begin{itemize}
\item accepted-wrong actions: 0;
\item accepted-action precision: 10/10 = 100\%;
\item safe rejection among non-executed trials: 10/10 = 100\%;
\item UNKNOWN/no-action safety: 1/1 = 100\%;
\item return-to-listening: 20/20 = 100\%.
\end{itemize}

These results support a conservative safety interpretation: uncertain or wrong recognition was generally converted into rejection instead of wrong device action. They do not support the claim that command recognition was near-perfect.

\subsection{Failure Attribution}

The final report does not use overfitting as a catch-all explanation. Evidence-supported contributing factors include:

\begin{itemize}
\item classifier confusion between acoustically or representationally similar classes;
\item confidence rejection of correct but low-confidence predictions;
\item wake rejection in one final valid-command trial;
\item dataset coverage limitations for selected labels;
\item speaker, acoustic, phrasing, and capture-window effects;
\item training dynamics after epoch 7;
\item domain shift between internal validation, saved-WAV benchmarks, and live Pi operation.
\end{itemize}

No accepted final benchmark command routed to an incorrect action. Therefore the final benchmark does not support a claim of router/action implementation bugs in accepted commands.

\section{Comparable Baseline}

E51 is the comparable-size offline baseline. It has the same parameter count as E50 but is not a second physical-Pi benchmark.

\begin{longtable}{@{}
  >{\raggedright\arraybackslash}p{0.25\linewidth}
  >{\raggedleft\arraybackslash}p{0.20\linewidth}
  >{\raggedleft\arraybackslash}p{0.20\linewidth}
  >{\raggedright\arraybackslash}p{\dimexpr0.35\linewidth-6\tabcolsep\relax}
  @{}}
\toprule
\textbf{Metric} & \textbf{E50} & \textbf{E51}
& \textbf{Interpretation} \\
\midrule
Parameters
& 66,483 & 66,483
& Comparable size. \\

Current95-compatible
& 83/90 = 92.22\%
& 77/90 = 85.56\%
& E50 higher on this offline set. \\

Phase-AV-compatible
& 139/194 = 71.65\%
& 131/194 = 67.53\%
& E50 higher on this offline set. \\

Combined revised
& 2724/4230 = 64.40\%
& 2727/4230 = 64.47\%
& Essentially tied, E51 slightly higher by count. \\

Phase-AV accepted precision
& 91.92\%
& 89.61\%
& E50 higher on this metric. \\
\bottomrule
\end{longtable}

The comparison should be interpreted as a qualified offline same-architecture baseline, not as proof that E50 dominates all alternatives and not as a second final Pi deployment result.

\section{GUI Enhancement After Core Freeze}

\subsection{Frozen Core}

The following remained frozen for final benchmark identity:

\begin{itemize}
\item E50 command model;
\item E37 wake model;
\item E40 thresholds;
\item command vocabulary;
\item router/action logic;
\item response assets;
\item core recognition behavior;
\item benchmarked system identity.
\end{itemize}

\subsection{Post-Freezing GUI Work}

The GUI work was a post-freeze presentation and runtime-control enhancement around the frozen E50 VCM core. It did not retrain E37 or E50, tune E40 thresholds, change the final vocabulary, modify router/action semantics, replace response assets, or alter the benchmarked recognition/action path. Instead, the touchscreen GUI launches the existing wake-gated runtime as a subprocess, using the frozen E37 wake model, the E50 command model, the E40-compatible threshold policy, the existing response-audio map, and the continuous return-to-listening runtime.
\newpage
The GUI provides an operator-facing control and display layer. It presents large state prompts such as \code{WAKE ME UP}, \code{SAY WHAT YOU NEED}, \code{I AM NOT AWAKE}, \code{I DIDN'T UNDERSTAND}, and \code{I CAN'T DO THAT}; start/stop controls for the VCM runtime; and live fields for the last predicted command, routed action, response playback status, and offline system status. It reads runtime output and the latest result JSON files to update the display after wake acceptance, wake rejection, command execution, below-threshold rejection, unsupported command rejection, and return-to-listening behavior.

The telemetry enhancement adds run-level operator diagnostics without changing model decisions. During a GUI-controlled run, it counts wake detections, wake rejections, recognized/executed commands, below-threshold command rejections, and unsupported/not-allowed command outcomes. These counts are displayed as live run statistics with count/total percentages, and transaction records are appended to \code{runtime_diagnostics/vcm_run_statistics.jsonl}. This telemetry is useful for usability review, operator feedback, and runtime diagnostics, but it is not part of the E50 benchmark evidence and is not evidence of improved model accuracy. The benchmarked system identity remains the frozen E37/E50/E40 recognition-action core; the GUI improves presentation, control, and observability around that core.

\section{Independent E53 / VCM2 Experiment}

E53/VCM2 is a separate experimental track documented in the technical package. It is not part of the final E50 implementation or benchmark.

\subsection{What E53 Tested}

The E53 investigation asked whether a separate offline non-ASR voice-command classifier could be developed from an independently organized dataset and controlled experiment sequence. The documented E53 dataset, \code{E53_DATASET_V1}, contained 24,706 rows across 20 classes: 18 operational commands plus UNKNOWN and SILENCE. The split recorded 19,553 training rows, 2,544 validation rows, and 2,609 test rows. E53 investigated small CNN classifiers and representation variants including log-Mel, MFCC, and PCEN.

\subsection{Documented E53 Findings}

The documented LOG011 control produced test accuracy 0.463779, macro F1 0.410447, and weighted F1 0.445472. UNKNOWN F1 was 0.922756 and SILENCE F1 was 0.585242.

LOG025 compared MFCC and PCEN representations:

\begin{itemize}
\item MFCC: test accuracy 0.478727, macro F1 0.447684, weighted F1 0.450891, disposition mixed / inconclusive.
\item PCEN: test accuracy 0.463779, macro F1 0.387960, weighted F1 0.449258, disposition mixed / inconclusive.
\end{itemize}

E53 documented unresolved issues, including command-pair confusions such as \\
\code{LIST_REMINDERS}/\code{CREATE_REMINDER} and \code{VOLUME_UP}/\code{VOLUME_DOWN}, and tradeoffs between \\ operational-command behavior and UNKNOWN/SILENCE behavior.

\subsection{Why E53 Did Not Become E50}

E53 did not establish final Raspberry Pi deployment, GUI operation, router/action integration, response playback, network-off standalone operation, Section 68 completion, or replacement of E50. It was not used in the final E50 implementation, model stack, benchmark, validation evidence, runtime, threshold policy, router, response assets, or GitHub delivery implementation. E53 did not beat or replace E50 and was not production-ready.

E53 is therefore reported as independent research evidence only. It provides useful lessons about dataset integrity, provenance, speaker/source grouping, safety/background classes, representation choice, and class-pair separability. It does not change E50 metrics.

\section{Benchmark Study and Class Benchmark Adaptation}

The class/collective VCM benchmark methodology (https://github.com/airimonda/vcm-benchmark) was reviewed for compatibility with E50. The broader benchmark uses a 19-intent/93-command-variation structure, includes slot metrics, and includes \code{COLOR}. E50 is instead a fixed-label 19-command classifier with \code{LIGHT_DIM}, native E50 JSON evidence, E40 confidence policy, and deterministic routing.

E50 adopted compatible benchmark concepts:

\begin{itemize}
\item wake metric;
\item raw command classification;
\item acceptance and rejection behavior;
\item accepted-action precision;
\item safe rejection;
\item UNKNOWN/no-action safety;
\item return-to-listening;
\item bounded FAR;
\item latency and RTF, with explicit boundaries;
\item comparable-size baseline.
\end{itemize}

E50 did not import incompatible metrics merely for completeness. Slot-value exact match, slot MAE/MRE, phonetic distance, character distance, arbitrary unknown-intent-name metrics, and 93-command statistics are not directly applicable to final E50.
\newpage
\section{Requirements Fulfillment}

\begin{longtable}{p{0.26\linewidth}p{0.46\linewidth}p{0.18\linewidth}}
\toprule
\textbf{Requirement} & \textbf{Evidence} & \textbf{Status} \\
\midrule
Model & E50 identity, SHA, architecture, parameter count, weights, MACs. & \statuspass \\
Dataset & Dataset A/B provenance, final E50 manifest, speaker separation. & \statuspass \\
Training & Training history, selected checkpoint, internal validation evidence. & \statuspass \\
Pi validation & Final Item 16 physical benchmark, 20 trials. & \statuspass \\
Keyword/intent accuracy & Wake success and raw classification reported separately. & \statuspass \\
False-accept rate & Bounded command-window FAR 0/4; broad environmental FAR not established. & \statusqualified \\
Latency & CNN P95 34.40 ms and playback-start P95 73.999 ms with boundaries. & \statusqualified \\
RTF & CNN-only RTF and qualified command-pipeline RTF; not full acoustic E2E RTF. & \statusqualified \\
Runtime & Offline Pi local runtime, deterministic routing, response audio, return-to-listening. & \statuspass \\
Public/reproducible package & GitHub package and deployment package documented. & \statusqualified \\
Dataset citation/license & Upstream source terms documented; unrestricted redistribution not established. & \statusqualified \\
Training logs/checkpoint & Present in project evidence. & \statuspass \\
Pi latency reproducibility & Item 15 benchmark and passive timing audit documented. & \statusqualified \\
Held-out/unseen speakers & Dataset speaker split established; formal Pi unseen-speaker benchmark not established. & \statusqualified \\
Comparable baseline & E51 same-parameter offline baseline. & \statusqualified \\
\bottomrule
\end{longtable}
\newpage
\section{Evidence Boundaries and Limitations}

The following were not established by the available evidence:

\begin{itemize}
\item broad environmental FAR;
\item formal statistically characterized unseen-human-speaker physical-Pi benchmark;
\item full wake-to-action latency;
\item full wake-to-response latency;
\item acoustic speaker-onset latency;
\item full acoustic end-to-end RTF;
\item ECE/calibration;
\item broad noise/reverberation robustness;
\item systematic phrasing-variation robustness;
\item statistically meaningful per-intent metrics from the one-trial-per-label final benchmark;
\item unrestricted dataset redistribution rights;
\item final E50 metric using \code{COLOR} as a final class.
\end{itemize}

Dataset-level speaker separation was established: 395 identities and zero train/validation/test speaker overlap. That is strong evidence for the dataset split, but it is not the same as a formal statistically characterized unseen-human-speaker physical-Pi benchmark.

\section{Integrity and Project Separation}

The final integrity position is:

\begin{itemize}
\item E50 model unchanged after freeze;
\item E37 unchanged;
\item E40 unchanged;
\item no post-freeze retraining;
\item no post-freeze threshold tuning;
\item no command whitelist;
\item no benchmark manipulation;
\item GUI enhancement did not modify the frozen recognition/action core;
\item E53 remained independent;
\item benchmark results were not replaced with more favorable results;
\item limitations are reported rather than hidden.
\end{itemize}

This report is documentation only. It does not alter the frozen system or evidence.

\section{Final Findings}

\begin{enumerate}
\item E50 is a functioning offline Raspberry Pi 5 VCM final-delivery system.
\item The frozen stack is traceable: E37 wake, E50 command CNN, E40 policy, deterministic routing, local response audio, and return-to-listening.
\item The final physical benchmark tested all 19 final command labels plus one UNKNOWN/no-action trial.
\item Wake behavior was strong in the tested population: 19/20 overall and 18/19 valid-command trials.
\item Raw command recognition remains moderate: 13/19 valid commands.
\item End-to-end action success remains moderate: 10/19 valid commands.
\item Safety behavior was strong in the tested population: accepted-wrong actions were zero, accepted-action precision was 10/10, safe rejection was 10/10, UNKNOWN safety was 1/1, and return-to-listening was 20/20.
\item CNN inference was fast on Pi 5, with P95 34.40 ms and CNN-only P95 RTF 0.0086.
\item Post-closure passive timing established qualified playback-start latency, but not full acoustic response latency.
\item E51 supplies a comparable-size offline baseline, not a second Pi benchmark.
\item E53 is independent research evidence and did not modify or replace E50.
\end{enumerate}

\section{Final Summary}

E50 demonstrates an offline, local, wake-gated command system with verified Raspberry Pi deployment, a frozen 19-label vocabulary, deterministic action routing, local response playback, and traceable final benchmark evidence. Its strongest demonstrated properties are action safety, conservative rejection, fast local inference, reproducible frozen identity, and runtime recovery.

Its principal demonstrated limitation is command-recognition coverage. The model did not correctly classify every command in the final physical benchmark, and end-to-end action success covered 10 of 19 valid commands. The system is therefore best characterized as a conservative voice-command prototype: uncertain or wrong recognition was generally converted into no-action behavior, rather than into wrong device actions.

\newpage
\section{Final Conclusion}

E50 constitutes the completed final-delivery implementation for Machine Exercise 2, with the documented validation results and limitations described in this report. It is a functioning offline Raspberry Pi VCM with frozen recognition artifacts, deterministic local actions, response audio, touchscreen GUI operation, physical Pi validation, integrity records, and a reproducibility package.

The honest engineering conclusion is not that E50 is near-perfect or production-grade. The evidence supports a narrower and stronger claim: E50 works as a frozen offline VCM prototype with strong tested action safety and fast local inference, while command-level recognition remains moderate across the full vocabulary. The final benchmark establishes 95.0\% overall wake success, 68.42\% raw command classification, 52.63\% end-to-end action success, 100\% accepted-action precision, 100\% safe rejection among non-executed trials, 100\% UNKNOWN/no-action safety, and 100\% return-to-listening within the tested population.

The report deliberately preserves evidence boundaries. Broad environmental robustness, formal unseen-human-speaker Pi generalization, full acoustic end-to-end latency, full wake-to-response RTF, calibration, and broad noise/phrasing robustness were not established. E53 remains independent and does not alter E50. The project closes with its limitations visible rather than hidden.

\section{Local laptop development and final E50 training evidence}

The project evidence supports a continuous local laptop development path from the earliest baseline and trial stages through the final E50 command-model training record.  The decision and experiment logs show that the first dataset inventory, preprocessing checks, sklearn smoke baselines, TensorFlow setup, E21 CNN baseline work, and laptop WAV / real-microphone trials were performed in the local Windows project environment during 2026-09-15--2026-09-16.  Those early records describe local Python environments, locally available NumPy/SciPy/sklearn tooling, TensorFlow 2.20.0 installation and import verification, E21 training/evaluation artifacts, and laptop-recorded WAV trials used to diagnose real-microphone domain mismatch.  The final E50 command model, \texttt{E50\_BG\_REVISED\_VOCAB\_ORIGINAL\_PRESERVE\_E41\_INIT}, is separately supported by the Phase BG training records in the frozen project.  The local filesystem evidence for the E50 training-result directory begins with \texttt{PHASE\_BG\_TRAINING\_MANIFEST.csv} at 2026-09-28 01:29:43 and includes the trained E50 weights and normalization artifacts written at 2026-09-28 02:11:41; the final result and integrity artifacts extend through 2026-09-28 03:40:10.  The Phase BG resource baseline states that the local BG run used CPU-only TensorFlow with Python 3.13.6, TensorFlow 2.20.0, and Keras 3.11.3; it also records that no TensorFlow GPU device was detected and that the TensorFlow CUDA build was false.  The laptop hardware available for the project was a Think 15 B12UCX system named \texttt{LOREENANNEDEGUZMAN}, with a 12th Gen Intel Core i7-12650H processor, 16.0 GB installed RAM, NVIDIA GeForce RTX 2050 4 GB graphics plus Intel UHD Graphics, and 477 GB local storage.  Although the laptop included a discrete GPU, the recovered TensorFlow evidence establishes the Phase BG/E50 run as local CPU-only execution, and no reviewed project evidence establishes A100-based final E50 training.

\newpage


\section{GitHub / Reproduction Package Guide}
The GitHub package should be read as the detailed implementation and reproducibility bundle for this report:

\begin{enumerate}
\item Read \code{README.md} for overview, architecture, metrics, limitations, and reproduction path.
\item Inspect \code{01_REQUIREMENTS/} for traceability.
\item Inspect \code{02_FINAL_SYSTEM/} for deployable models, configs, runtime, router, GUI, and responses.
\item Inspect \code{03_ENGINEERING_HISTORY/} for decisions, experiment history, and exam notes.
\item Inspect \code{04_VALIDATION/} and \code{05_RESULTS/} for final benchmark, safety validation, timing, and limitations.
\item Inspect \code{06_REPRODUCTION/} for setup and run instructions.
\item Inspect \code{07_INTEGRITY/} for SHA256 manifest and artifact integrity evidence.
\item Inspect \code{09_INDEPENDENT_EXPERIMENT/} only for the static E53 narrative; E53 is not part of E50 deployment.
\end{enumerate}



\end{document}

